% \section{Problem reformulation}
% \label{Section: Problem Reformulation}
% \subsection{W-type Regularized Deep Hashing Model}
% Deep hashing methods aim to learn network parameters \(\bx \in \RR^{d_{\bx}}\) to generate discrete codes that preserve semantic similarity effectively in the Hamming space. 
% A natural approach is to minimize a similarity-preserving loss function \( \frac{1}{n}\sum_{i=1}^n L_s(\text{sgn}(\bD_i(\bx)))\) where $n$ is the number of training samples, $\bD_i(\bx)$ represents the \(i\)-th network output and \(f_i\) is the similarity-preserving loss function (e.g., pairwise cross-entropy loss \cite{zhuDeepHashingNetwork2016}). However, the function \( sgn(\cdot) \) renders standard backpropagation infeasible due to its zero gradient for all nonzero inputs. This discreteness poses a central challenge in deep hashing.

% A W-type regularization (e.g., $\lambda | |z| - 1 |$) has effectively mitigated this issue. This type of regularization encourages outputs to approach binary values \(\{-1, 1\}\) during training. Specifically, the model is trained to produce continuous outputs \(\bu_i = \bD_i(\bx)\) with W-type regularization, and binary codes are obtained post-training via \(\bb_i = \text{sgn}(\bu_i)\).  Consequently, we formulate the core optimization problem for W-type regularized deep hashing methods as follows,
% \begin{align}
% \label{Eq: the primal deep hashing problem}
% \min_{\bx \in \RR^{d_{\bx}}} \quad \frac{1}{n} \sumn f_i (\bx) + \sumn h(\bD_i(\bx)), 
% \end{align}where \(\frac{1}{n}\sumn f_i\) is the similarity-preserving loss, \(D_i\) is the \(i\)-th network output and \(h\) is a W-type regularization.
% % \footnote{In the problem \eqref{Eq: the primal deep hashing problem}, we use vectorized \(\bz\), thus the regularization is also \(\| |\bz| - \bm{1}\|_1 = \sumn | |z_i| - 1 |\)}.
% \subsection{A Two-Block Structured Reformulation}
% We then employ variable splitting % \cite{wangLiftedBregmanTraining2022, zengBlockCoordinateDescent2018, lauProximalBlockCoordinate2018}, ]\
% to reformulate problem \eqref{Eq: the primal deep hashing problem} equivalently as the constrained problem:
% \begin{align} 
% \label{Eq: constrained_dlh} 
% \min_{\bx \in \RR^{d_{\bx}}, \bB \in \RR^{n\times d}}&\quad \frac{1}{n} \sumn f_i (\bx) + \sum_{i=1}^n h(\bb_i) \notag \\ 
% \text{s.t.} &\quad \bb_i = \bD_i(\bx), \; \forall i=1,\dots,n. 
% \end{align} Combined with the quadratic penalty method \cite{nocedal1999numerical}, this transformation yields the following two-block finite-sum optimization problem:
% \begin{align} 
% \label{Prop: block-wise regularized problem} 
% \nonumber&\min_{\bx \in \RR^{d_{\bx}}, \bB \in \RR^{n\times d}} J(\bx, \bB) \\
% &= \overbrace{\underbrace{\frac{1}{n}\sum_{i=1}^n f_i(\bx)}_{f(\bx)} + \underbrace{\frac{\gamma}{2n} \sum_{i=1}^n \| \bD_i(\bx) - \bb_i\|^2}_{P(\bx, \bB)}}^{F(\bx, \bB) = \frac{1}{n}\sumn F_i(\bx, \bb_i)} + \underbrace{\sum_{i=1}^n h(\bb_i)}_{h(\bB)}. \tag{P}
% \end{align} Here, \( \bB = [\bb_1, \bb_2, \dots, \bb_n]\) and \( \gamma > 0 \) is the penalty parameter. This reformulation decomposes the original problem \eqref{Eq: the primal deep hashing problem} into simpler structural blocks, enabling us to exploit the problem structure better. Building upon this problem, we develop a stochastic primal-dual algorithm that effectively addresses the challenges above.
\section{Problem Reformulation}
\label{Section: Problem Reformulation}
\subsection{W-type Regularization in Deep Hashing Model}
\label{Sec:W-type Regularized Deep Hashing Model}
Deep hashing aims to learn network parameters \(\bx \in \RR^{\blue{q}}\) that generate binary codes preserving semantic similarity in the Hamming space. A direct formulation is
\begin{align}
\label{Eq:direct_deep_hashing_problem}
\blue{\min_{\bx \in \RR^{q}}\ \frac{1}{n}\sum_{i=1}^{n}\ell_{s,i}\!\big(\operatorname{sgn}(\mathcal{D}(\bx))\big),}
\end{align}
where \(n\) is the number of training samples, \blue{\(\bD_i:\RR^q\to\RR^d\) is the continuous network output mapping for sample \(i\), and \(\ell_{s,i}:\RR^{n\times d}\to\RR\) is a similarity-preserving loss} (e.g., pairwise cross-entropy \cite{zhuDeepHashingNetwork2016}). This objective is hard to optimize by backpropagation because \(\operatorname{sgn}(\cdot)\) is nondifferentiable and has zero gradient almost everywhere. 

\blue{To address this issue, penalties whose graphs are W-shaped, vanishing exactly at
$\pm1$, are added to drive continuous outputs toward $\{-1,1\}^d$. They are standard in
deep hashing~\cite{liuDeepSupervisedHashing,wang2016DTSH,zhuDeepHashingNetwork2016,
yuanCentralSimilarityQuantization2020,he2025mambahash,yin2026crh}.
Related distance-based regularizers are also used to train binary and quantized
networks~\cite{yin2018binaryrelax,bai2019proxquant}. We formalize the scalar penalties
considered here as the W-type class in \Cref{def:wtype}.}

\begingroup
\color{blue}
\begin{definition}[W-type regularizer]\label{def:wtype}
Let \(\phi:\mathbb{R}\to[0,+\infty)\). For \(\bz\in\mathbb{R}^{d}\) and \(\bZ\in\mathbb{R}^{n\times d}\), define \(h(\bz):=\sum_{j=1}^{d}\phi(z_j)\) and \(h(\bZ):=\sum_{i=1}^{n}\sum_{j=1}^{d}\phi(z_{ij})\). We call \(h\) a \emph{W-type regularizer} if \(\phi\) satisfies:
\begin{enumerate}[label=\textup{(W\arabic*)}]
  \item\label{W1} \(\phi\) is even and continuous, and \(\phi(t)=0\) if and only if \(t\in\{-1,1\}\);
  \item\label{W2} \(\phi\) is convex on \([1,+\infty)\);
  \item\label{W3} \(\phi\) decreases at a nonvanishing rate on \([0,1)\): for every \(b\in(0,1)\) there exists \(c_b>0\) such that \(\phi(s)-\phi(t)\ge c_b(t-s)\) for all \(0\le s<t\le b\).
\end{enumerate}
\end{definition}



Each condition is used once: \ref{W1} makes the zero set of \(h\) exactly \(\{-1,1\}^{n\times d}\); \ref{W2} yields the envelope formula of Lemma~\ref{lem:wtype-conj}, so that the dual reformulation leaves \(h\) unchanged on \(|t|\ge1\); and \ref{W3} excludes stationary points in \((-1,1)\) (Lemma~\ref{lem:no-inner-stationary}), which gives the stationarity-to-contact implication of Proposition~\ref{prop:transfer}.

\begin{table}[t]
\centering
\color{blue}
\caption{Representative W-type regularizers, with \(\lambda,\delta>0\), \(\psi(u):=\frac12[(1+u)\log(1+u)+(1-u)\log(1-u)]\), and \(\operatorname{Hub}_{\delta}(s):=s^2/(2\delta)\) if \(|s|\le\delta\), \(|s|-\delta/2\) otherwise. Outside \(\operatorname{dom}\phi^*\), \(\phi^*=+\infty\); \(\operatorname{clip}\) denotes projection onto \([-\lambda,\lambda]\); ``---'' means that no closed form is available.}
\label{tab:wtype}
\resizebox{\linewidth}{!}{
\begin{tabular}{lllcccc}
\toprule
\multicolumn{2}{l}{\(\phi(t)\)} & Used in & \(\partial\phi(0)\) & \(\phi^*(y)\) & \(\operatorname{dom}\phi^*\) & \(\operatorname{prox}_{\sigma\phi^*}(v)\)\\
\midrule
\multirow{2}{*}{\(\lambda\bigl||t|-1\bigr|^{p}\)}
 & \(p=1\) & \cite{liuDeepSupervisedHashing,wang2023deep,bai2019proxquant}
 & \(\{\pm\lambda\}\) & \(|y|\) & \([-\lambda,\lambda]\)
 & \(\operatorname{clip}\bigl(\operatorname{sgn}(v)(|v|-\sigma)_+\bigr)\)\\
 & \(p=2\) & \cite{wang2016DTSH,he2025mambahash,semanticHashCenters2025,yin2026crh,bai2019proxquant}
 & \(\{\pm2\lambda\}\) & \(|y|+\frac{y^2}{4\lambda}\) & \(\mathbb{R}\)
 & \(\operatorname{sgn}(v)\frac{(|v|-\sigma)_+}{1+\sigma/(2\lambda)}\)\\
\midrule
\multicolumn{2}{l}{\(\lambda\log\cosh(|t|-1)\)} & \cite{zhuDeepHashingNetwork2016,yuanCentralSimilarityQuantization2020}
 & \(\{\pm\lambda\tanh1\}\)
 & \(|y|+\lambda\psi(|y|/\lambda)\)
 & \([-\lambda,\lambda]\) & ---\\
\multicolumn{2}{l}{\(\lambda\operatorname{Hub}_{\delta}(|t|-1)\)} & ---
 & \(\{\pm\lambda\min\{1,\frac1\delta\}\}\) & \(|y|+\frac{\delta y^2}{2\lambda}\) & \([-\lambda,\lambda]\)
 & \(\operatorname{clip}\Bigl(\operatorname{sgn}(v)\frac{(|v|-\sigma)_+}{1+\sigma\delta/\lambda}\Bigr)\)\\
\bottomrule
\end{tabular}}
\end{table}

All rows of Table~\ref{tab:wtype} satisfy \ref{W1}--\ref{W3}, and none is differentiable at \(0\): its subdifferential there is \(\{\pm c\}\) with \(c>0\). This is forced by the class itself, since an even \(\phi\) that is differentiable at \(0\) has \(\phi'(0)=0\), contradicting \ref{W3}. The rows differ in where else they are nonsmooth and how fast they grow. For \(p=1\) the kinks also lie at \(\pm1\), where the binary codes are. The other three rows are smooth at \(\pm1\), but the kink at \(0\) still rules out \(L\)-smoothness and weak convexity, so standard complexity guarantees for gradient methods do not apply to them either. Tail growth determines \(\operatorname{dom}\phi^*\): smoothing \(p=1\) by squaring gives \(\operatorname{dom}\phi^*=\mathbb{R}\), whereas smoothing it by \(\log\cosh\) or \(\operatorname{Hub}_\delta\) keeps \(\operatorname{dom}\phi^*=[-\lambda,\lambda]\).

\setcounter{lemma}{1}
\begin{lemma}\label{lem:wtype-conj}
Let \(h\) be a W-type regularizer. Then \(\phi^*\) is proper, lower semicontinuous, convex, and even, \(\phi^*(y)\ge|y|\) for all \(y\), and \(\operatorname{dom}\phi^*\) has nonempty interior. Moreover,
\[
  \phi^{**}(t)=\phi\bigl(\max\{|t|,1\}\bigr),\qquad
  \phi^{**}(t)=\phi(t)\iff |t|\ge1.
\]
Consequently, \(h^*(\bLmbd)=\sum_{i,j}\phi^*(\Lambda_{ij})\), \(h^{**}(\bB)=\sum_{i,j}\phi^{**}(b_{ij})\), and \(h(\bB)=h^{**}(\bB)\) if and only if \(|b_{ij}|\ge1\) for all \(i,j\).
\end{lemma}
\begin{proof}
Let \(\psi(t):=\phi(\max\{|t|,1\})\). By \ref{W1}--\ref{W2}, \(\phi\) is convex and nonnegative on \([1,+\infty)\), with \(\phi(1)=0\); hence it is nondecreasing there. Thus \(s\mapsto\phi(\max\{s,1\})\) is convex and nondecreasing on \([0,+\infty)\), so \(\psi\) is convex and satisfies \(\psi\le\phi\). Conversely, any convex \(q\le\phi\) satisfies \(q(\pm1)\le0\), hence \(q\le0=\psi\) on \([-1,1]\), while \(q\le\phi=\psi\) on \(|t|\ge1\). Therefore \(\phi^{**}=\psi\), the largest lower semicontinuous convex minorant of \(\phi\)~\cite[Theorem~12.2]{rockafellar1970convex}. For \(|t|<1\), \ref{W1} gives \(\phi(t)>0=\psi(t)\), proving the contact-set formula.

Evaluating the definition of \(\phi^*\) at \(t=\pm1\) gives \(\phi^*(y)\ge|y|\). Since \(\phi(t_0)>0\) for any \(t_0>1\), convexity on \([1,+\infty)\) implies at least linear growth for large \(t\); evenness gives the same for negative \(t\). Hence \(\operatorname{dom}\phi^*\) contains an open interval around zero. The separable formulas follow from~\cite[Theorem~16.4]{rockafellar1970convex}.
\end{proof}

\begin{remark}[Scope of Definition~\ref{def:wtype}]\label{rem:wtype-scope}
Conditions \ref{W2} and \ref{W3} cannot be dropped.
\begin{enumerate}[label=\textup{(\roman*)}]
  \item \emph{Without \ref{W3}.} The function \(\phi(t)=|t^2-1|\) satisfies \ref{W1}--\ref{W2} but violates \ref{W3}, because its slope vanishes at \(0\); indeed \(\partial\phi(0)=\{0\}\). In the setting of Example~\ref{ex:gap}, its convex envelope is \(\phi^{**}(t)=(t^2-1)_+\), and \((b,\Lambda)=(0,0)\) satisfies the KKT conditions. It is also stationary for Problem~(P), but \(J(b)=1-\tfrac12b^2\) near zero, so it is a local maximizer. Here \(|b|<1\) and \(h(0)\ne h^{**}(0)\), showing why \ref{W3} is needed for Proposition~\ref{prop:transfer}.
  \item \emph{Without \ref{W2}.} The function \(\phi(t)=\lambda\log\bigl(1+\bigl||t|-1\bigr|\bigr)\) satisfies \ref{W1} and \ref{W3}, but its tails are concave. Its sublinear growth gives \(\operatorname{dom}\phi^*=\{0\}\) and \(h^{**}\equiv0\), so the dual reformulation forces \(\bLmbd=\zero\) and removes the regularizer.
\end{enumerate}
Finally, the row \(p=2\) of Table~\ref{tab:wtype} satisfies \ref{W1}--\ref{W3}, so the structural optimality propositions remain valid for it. Its conjugate has full domain, however, and the current complexity analysis, which uses boundedness of the dual domain, does not cover this quadratic-growth case. Condition~\ref{W1} specifies the binary-code zero set and is therefore part of the model definition.
\end{remark}
\endgroup
Specifically, the model is trained to produce the \(i\)-th continuous network output \(\bD_i(\bx)\) with W-type regularization \(h\), and binary codes are obtained post-training via \(\operatorname{sgn}(\bD_i(\bx))\blue{\in\{-1,1\}^d}\). \blue{Consequently, we formulate the following core optimization problem for W-type regularized deep hashing methods, which serves as a continuous surrogate of Problem~\eqref{Eq:direct_deep_hashing_problem}:}
\begin{align}
\label{Eq: the primal deep hashing problem}
\min_{\bx \in \RR^{\blue{q}}} \quad \frac{1}{n}\sum_{i=1}^n f_i(\bx) + \sum_{i=1}^n h(\bD_i(\bx)),
\tag{P0}
\end{align}\blue{where
\(f_i(\bx):=\ell_{s,i}\!\left(\mathcal{D}(\bx)\right)+\frac{\mu}{2}\|\bx\|^2\) with \(\mu>0\) ensures coercivity.}

We then employ variable splitting by introducing \blue{continuous auxiliary variables \(\bb_i\in\RR^d\) and \(\bB=[\bb_1^\top;\ldots;\bb_n^\top]\in\RR^{n\times d}\), subject to \(\bb_i=\bD_i(\bx)\) for all \(i=1,\dots,n\),} which yields the following equivalent nonlinearly constrained problem:
\begin{align}
    \blue{\min_{\bx \in \RR^{q},\ \bB \in \RR^{n\times d}}\quad \frac{1}{n}\sum_{i=1}^{n} f_i(\bx)+\sum_{i=1}^{n}h(\bb_i)\quad \text{s.t.}\quad \bb_i=\bD_i(\bx),\ \forall i=1,\dots,n.} \notag
\end{align}
We then \blue{relax these constraints by} the quadratic penalty method \cite[Sec.~17.1]{nocedalWright2006} to obtain the following two-block finite-sum optimization problem:
\begin{align}
    \label{Prop: block-wise regularized problem}
    \blue{\min_{\bx \in \RR^{q},\,\bB \in \RR^{n\times d}}\ J(\bx,\bB):=F(\bx,\bB)+h(\bB),} \tag{P}
\end{align}
\blue{Here, \(\gamma>0\) is the penalty parameter, and we use the separable extension \(h(\bB):=\sum_{i=1}^{n}h(\bb_i)\) for matrices. We define}
\begin{align}
    \blue{F(\bx,\bB)} &\blue{:= \frac{1}{n}\sum_{i=1}^{n}F_i(\bx,\bb_i),} \notag\\
    \blue{F_i(\bx,\bb_i)} &\blue{:= f_i(\bx)+\frac{\gamma}{2}\|\bD_i(\bx)-\bb_i\|^2,} \notag
\end{align}
\blue{Averaging over \(i\) gives the decomposition \(F=f+P\), where}
\begin{align}
    \blue{f(\bx)} &\blue{:= \frac{1}{n}\sum_{i=1}^{n}f_i(\bx),} \notag\\
    \blue{P(\bx,\bB)} &\blue{:= \frac{\gamma}{2n}\sum_{i=1}^{n}\|\bD_i(\bx)-\bb_i\|^2.} \notag
\end{align}
This reformulation decomposes the original problem \eqref{Eq: the primal deep hashing problem} into simpler structural blocks, enabling us to exploit the problem structure better. Building upon Problem~\eqref{Prop: block-wise regularized problem}, we develop a stochastic primal-dual algorithm {that addresses the lack of tractable proximal mappings in the nonsmooth composite form and the absence of convergence guarantees in existing deep hashing methods.}
